%  05_evaluation.tex
\section{Evaluation Protocol}
\label{sec:evaluation}

\subsection{Gallery and probe construction}

For each test split we designate one image per identity as the \emph{gallery} (the
enrolment reference) and treat the remaining images as \emph{probes}.  Gallery entries are
selected deterministically by sorting by file path and taking the first image; this ensures
reproducibility.  Identities with fewer than $g + 1$ images (gallery size $g=1$ plus at
least one probe) are excluded from evaluation.

\subsection{Verification}

Verification asks: \emph{do two images show the same individual?}
All $\binom{N_g + N_p}{2}$ unordered pairs are generated from the combined gallery and
probe pool; a pair is labelled positive if both images share the same identity.
The similarity score for each pair is the cosine similarity of L2-normalised embeddings:
\begin{equation}
  s(\mathbf{x}_i, \mathbf{x}_j) = \frac{\hat{\mathbf{z}}_i \cdot \hat{\mathbf{z}}_j}
                                         {\max(\|\hat{\mathbf{z}}_i\|, \varepsilon)
                                          \max(\|\hat{\mathbf{z}}_j\|, \varepsilon)},
  \quad \varepsilon = 10^{-12}.
  \label{eq:cosine}
\end{equation}
The following metrics are computed:

\paragraph{ROC-AUC.}  Area under the receiver operating characteristic curve, computed via
\texttt{sklearn\discretionary{.}{}{.}metrics\discretionary{.}{}{.}roc\_auc\_score}.
Provides a threshold-independent measure of discriminative power.

\paragraph{Youden's J and optimal threshold.}  $J = \mathrm{TPR} - \mathrm{FPR}$
(equivalently, sensitivity + specificity $- 1$)~\cite{youden1950}, maximised over the ROC
curve to yield an optimal decision threshold.  This is used for F\textsubscript{1},
precision, recall, and specificity computation.

\paragraph{Threshold-selection policy.}
All operating-point thresholds (Youden's J, $\mathrm{Recall@FAR}\le10^{-2}$, and the
TAR@FAR budgets) are selected on the \emph{validation} split and then applied
unchanged to the test split; the resulting test operating points are reported under
the \texttt{test\_val\_threshold/*} prefix.  Thresholds re-fit directly on the test
ROC curve are retained only as diagnostic upper bounds, never as headline numbers, to
avoid optimistic bias from tuning the decision boundary on the evaluation set.  The
permutation test below likewise uses the validation-derived threshold so that the null
distribution reflects a deployable, test-blind operating point.

\paragraph{TAR@FAR.}  The True Acceptance Rate at a fixed False Acceptance Rate (FAR) budget,
defined as the maximum TPR among ROC operating points with $\mathrm{FPR} \le
\mathrm{FAR}_{\text{target}}$.  We report TAR at $\mathrm{FAR} = 10^{-3}$ and $10^{-4}$.

\paragraph{Recall@FAR$\le 10^{-2}$ (headline metric).}
With test partitions of at most a few hundred pairs, a FAR budget of $10^{-3}$ permits
fewer than one false match and degenerates to noise.  We therefore adopt
$\mathrm{Recall@FAR} \le 10^{-2}$ as the primary biometric operating metric.  This
corresponds to the TAR at the most lenient operating point not exceeding a 1\% false
positive rate, offering a meaningful trade-off between security and usability at the
current dataset scale.

\subsection{Identification}

Identification asks: \emph{given a probe image, which gallery identity does it belong to?}
For each probe, cosine similarities to all gallery entries are ranked in descending order.
Identification metrics are:

\begin{itemize}[itemsep=2pt]
  \item \textbf{Top-1 / Top-5 accuracy}: fraction of probes for which the correct identity
        is ranked first (or within the top 5).
  \item \textbf{CMC curve}: cumulative match characteristic up to rank $k = 10$,
        $\mathrm{CMC}[k] = $ fraction of probes with first correct match at rank $\le k$.
  \item \textbf{mAP}: mean average precision in the information-retrieval sense, averaging
        the per-probe average precision (AP) over all probes.
\end{itemize}

\subsection{Statistical significance: permutation test}

To assess whether observed performance exceeds chance, we run an \emph{inference-mode
permutation test}: identity labels in the test manifest are permuted $n = 100$ times; for
each permutation, embeddings are re-computed (no retraining) and the target metric is
evaluated.  The one-sided p-value with continuity correction is:
\begin{equation}
  p = \frac{\#\{\mathrm{null} \ge \mathrm{observed}\} + 1}{n + 1}.
  \label{eq:pvalue}
\end{equation}
Significance threshold $\alpha = 0.05$.  The inference-mode variant is substantially faster
than full-retrain permutation testing while still providing a valid null distribution for
the embedding quality.
